%
%pdf 
%pdf structure
%
%author :Tariq Ali
%

\chapter{Pdf}

\section{Structure of Pdf}

PDF is a portable document format that can be used to present documents that include text, images, multimedia elements, web page links and more. PDF has more functions than just text, it can include images and other multimedia elements, be password protected. \cite{Lukan:2020}

\begin{figure}[h]
	\centering
	\includegraphics{Pdf/PDFStructure}
    \caption{PDF structure}
\end{figure}

Every PDF document has the following elements:

\subsection{Header}

This is the first line of a PDF file and specifies the version number of the used PDF specification which the document uses. If we want to find that out, we can use the hex editor or simply use the xxd command as below:

\begin{alltt}
	\%PDF-1.2
\end{alltt}

The temp.pdf PDF document uses the PDF specification 1.2. The ‘\%’ character is a comment in PDF, so the above example presents the first- and second line being comments, which is true for all PDF documents. Currently the version numbers are of the form 1.N, where the N is from range 0-7. \cite{Steve:2018}

\subsection{Body}

In the body of the PDF document, there are objects that typically include text streams, images, other multimedia elements, etc. The Body section is used to hold all the document’s data being shown to the user.

\subsection{Cross-Reference Table}

This table is similar to a directory. It contains the location of each object within the PDF file. By looking at the entries in this table, the PDF reading application can easily locate an object within the file. This saves time as the object is accessed in a random manner. The cross-reference table can have one or more sections. Each of these sections can have one or more subsections.\\
Each section begins with the word 'xref'. Following this line are two numbers separated by a single space. The first number is the object number of the first of the series of objects listed below it. The second number refers to the number of entries in that subsection.Each object is represented by one entry in the cross reference table, which is always 20 bytes long. For example:
\begin{alltt}
		xref
		0 1
		0000000023 65535 f
		3 1
		0000025324 00000 n
\end{alltt}

In the example above, we can see that we have two subsections (note the two lines that only contain two numbers). The first number in those lines corresponds to the object number, while the second line states the number of objects in the current subsection.

\begin{itemize}
	\item The ten digit value shows how far the object is from the start of the file, 23 denotes that the object is 23 bytes from the start of the file.
	\item The five digit number specifying the object generation number.
	\item 'n' denotes that the object is still in use and 'f' denotes that the object has been deleted and is free.
\end{itemize}
The first object has an ID 0 and always contains one entry with generation number 65535 that is at the head of the list of free objects.\\
\\
The second subsection has an object ID 3 and contains one element, the object 3 that starts at an offset 25324 bytes from the beginning of the document.

\subsection{Trailer}

The end of a PDF file is read first by the PDF reading application. The trailer holds information about the location and details of the Cross-reference table. The trailer has three parts.

\begin{itemize}
	\item The first part has the keyword trailer followed by a dictionary that holds values for certain fields.
	\item The second part has the keyword startxref, and in the next line, a number. The number denotes how far (in bytes) the keyword xref (of the last section of the cross-reference table) is from the start of the file.
	\item The very next line has the value \%\%EOF to denote the end of the file.
\end{itemize}
For example:

\begin{alltt}
	trailer
	<< /Size 62 /Root 1 0 R /Info 2 0 R
	/ID [(X$X@>66...)(X$X@>66...)]
	>>
	startxref
	361441
	\%\%EOF
\end{alltt}

Size is the total number of entries in the cross-reference table, in the example above it is 62 entries.

\section{Character set of PDFs}

There are basically three tyoes of character sets used in PDF.

\subsection{White-space character}

White-space characters separate names and other objects from each other.

\subsection{Delimiter characters}

(, ), <, >, [, ], {, }, / and \% (4 pairs and 2 unique).These are used in the objects. They basically act as delimiters to mark the boundary or border for entities.

\subsection{Regular characters}

All characters other than White-space and Delimiter characters.

\section{PDF Data types}

The PDF document contains eight basic types of objects, it includes

\subsection{Booleans}

The keyword \textbf{true} and \textbf{false} represents the boolean values.

\subsection{Number Data}

There are two types of numbers in a PDF document

\begin{enumerate}
	\item \textbf{Integar} consists of one or more digits can also be preceded by a plus or minus sign. For example:
	\begin{itemize}
		\item 123, +123, -123
	\end{itemize}
    \item A \textbf{real} value can be represented with one or more digits with an optional sign(+ or -) with a decimal point. For example:
    \begin{itemize}
    	\item 123.0, -123.0, +123.0, 123., -.123
    \end{itemize}
\end{enumerate}

\subsection{Names}

Names consist of a sequence of characters. A forward slash / must be used to introduce the name. In case hash (\#) is part of the name, use \# followed by its hexadecimal code 23. To represent characters using their hexadecimal value use \# followed by the hexadecimal value. All characters that are not regular characters have to be represented by the \# followed by their hexadecimal value.

\begin{alltt}
	/My\#20Name represents My Name
	/All\#20\#23Numbers represents All \#Numbers
\end{alltt}

The examples of names can be seen in the table \ref{fig:PDFNames}.

\begin{figure}
	\centering
	\includegraphics{Pdf/NamesObject}
	\caption{PDF Names} \label{fig:PDFNames}
\end{figure}

\subsection{Strings}

Strings contain characters. They can be literal characters within parenthesis or hexadecimal data within angle brackets.

\begin{itemize}
	\item A string embedded in parantheses can be represented as:

\begin{alltt}
	(mystring)
\end{alltt}

    \item A string with hexadecimal numbers may be represented as:

\begin{alltt}
	<48454C4C4F>
\end{alltt}
\end{itemize}

Each pair is taken as a value. In the above example, the hexadecimal value \#48 is decimal 72 which is the ASCII equivalent of H. Similarly \#45 is E, \#4C is L and \#4F is O. The above string is same as the string (HELLO).

\subsection{Arrays}

Arrays in PDF documents are represented as a sequence of PDF objects, which may be of different types and enclosed in square brackets. An array in a PDF can hold objects like strings, names, numbers, dictionaries and even other arrays.

\begin{verbatim}
	[false 170 85.5 (Hello) /My#20Name]
\end{verbatim}

\subsection{Dictionary}

These are similar to an actual dictionary, where a description follows a word. The description here can be any object (including another dictionary) but the name is always a Name object. The key (Name) is unique (there cannot be two similar Names). A dictionary is represented within << >>. For example:

\begin{alltt}
	<< /Name (Steve)
	/Age 99
	/Address << /Number 1234
	                /Street (Java Street)
	                /Suburb (Java Town)
	                /PostCode (4321)
	>>
	>> 
\end{alltt}

\subsection{Streams}

A stream object is represented by a sequence of bytes and may be unlimited in length, which is why images and other big data blocks are usually represented as streams. The contents of each page in PDF is represented as a 'Contents' stream. It consists of a dictionary followed by the keyword ‘stream’, newline, the stream’s data and the keyword ‘endstream’. For example:

\begin{alltt}
	dictionary
	stream
	……….
	endstream
\end{alltt}

The stream consists of the data, the dictionary contains information about the stream itself. Mostly stream dictionaries have common keys out of which the only mandatory key is Length.

\begin{itemize}
	\item \textbf{Length - Mandatory entry} that contains the length (number of bytes) of the stream. An error occurs if the stream has more bytes of data that the length mentioned in the dictionary.
\end{itemize}

\subsection{Null Object}

The null object is represented by a keyword “null”



